\documentclass[border=10pt]{standalone}
\usepackage{tikz}
\usepackage[expansion=false]{microtype}
\usetikzlibrary{arrows.meta, positioning, fit, backgrounds, shapes.geometric}
\usepackage{xcolor}

\definecolor{abstrcol}{HTML}{2C5F8A}
\definecolor{kaldocol}{HTML}{8B3A3A}
\definecolor{phonocol}{HTML}{2D6A4F}
\definecolor{paramcol}{HTML}{6B7280}

\begin{document}

\begin{tikzpicture}[
  >=Latex,
  font=\sffamily\small,
]

% =========================================================
% GRID PARAMETERS
% =========================================================
% Five columns at x = 0, 4.6, 9.2, 13.8, 18.4
% Three rows at y =  0 (abstract), -3.0 (kaldo), -5.6 (phonopy)
%   plus a sub-row at y = -1.8 (FC order=3, abstract only)
% =========================================================

\def\colA{0}
\def\colB{4.6}
\def\colC{9.2}
\def\colD{13.8}
\def\colE{18.4}

\def\rowABS{0}
\def\rowABS3{-1.8}
\def\rowKAL{-3.6}
\def\rowKAL3{-5.0}
\def\rowPHO{-6.4}
\def\rowPHO3{-7.8}

% Styles --------------------------------------------------
\tikzset{
  abstr/.style    = {draw, rounded corners=2pt, minimum width=3.4cm, minimum height=1.0cm,
                     fill=abstrcol!12, draw=abstrcol, line width=0.5pt, align=center},
  param/.style    = {draw, rounded corners=2pt, minimum width=3.4cm, minimum height=1.0cm,
                     fill=paramcol!12, draw=paramcol, line width=0.5pt, dashed,
                     align=center, font=\sffamily\itshape},
  matnode/.style  = {draw, rounded corners=2pt, minimum width=3.8cm, minimum height=1.3cm,
                     line width=0.5pt, align=left, text width=3.6cm, inner sep=4pt,
                     font=\scriptsize},
  kaldo/.style    = {matnode, fill=kaldocol!6, draw=kaldocol},
  phono/.style    = {matnode, fill=phonocol!6, draw=phonocol},
  oparr/.style    = {->, thick, abstrcol},
  stubarr/.style  = {->, thick, dashed, paramcol},
  matarr/.style   = {-, dotted, gray!70, line width=0.4pt},
  oplabel/.style  = {font=\scriptsize, midway, fill=white, inner sep=1.5pt,
                     text=abstrcol!80!black},
  rowlbl/.style   = {anchor=east, font=\sffamily\bfseries},
}

% =========================================================
% ABSTRACT DAG (top row)
% =========================================================

\node[param] (POT)  at (\colA,\rowABS)  {Potential};
\node[abstr] (FC2)  at (\colB,\rowABS)  {ForceConstants\\\footnotesize{[order=2]}};
\node[abstr] (DM)   at (\colC,\rowABS)  {DynamicalMatrix};
\node[abstr] (DISP) at (\colD,\rowABS)  {Dispersion};

% FC order=3 branch (sub-row, abstract only — its materializations slot beneath FC2)
\node[abstr] (FC3) at (\colB,\rowABS3) {ForceConstants\\\footnotesize{[order=3]}};

% Operations
\draw[stubarr] (POT) -- node[oplabel, above] {\textbf{extract+trunc}}
                       node[oplabel, below, font=\tiny] {order=2,\,method=$M$} (FC2);
\draw[stubarr] (POT.south east) to[bend right=15]
                       node[oplabel, sloped, below, font=\tiny, pos=0.55] {order=3,\,method=$M'$}
                       (FC3.west);
\draw[oparr] (FC2) -- node[oplabel, above] {\textbf{Fourier}}
                     node[oplabel, below, font=\tiny] {q-mesh} (DM);
\draw[oparr] (DM)  -- node[oplabel, above] {\textbf{diag}} (DISP);

% Row label
\node[rowlbl, color=abstrcol]
     at ([xshift=-0.4cm]POT.west |- POT) {abstract};

% =========================================================
% KALDO ROW
% =========================================================

\node[kaldo] (POTk) at (\colA,\rowKAL)
     {\textbf{kaldo (Tersoff)}\\
      \texttt{tersoff(Si.tersoff)}\\
      \emph{opaque label}};

\node[kaldo] (FC2k) at (\colB,\rowKAL)
     {\textbf{kaldo}\\
      \texttt{ForceConstants.from\_folder(}\\
      \quad\texttt{format='lammps',}\\
      \quad\texttt{supercell=(3,3,3))}\\
      file: \texttt{Dyn.form}};

\node[kaldo] (DMk) at (\colC,\rowKAL)
     {\textbf{kaldo}\\
      \texttt{phonons.dynmat} (lazy)\\
      module:\\
      \texttt{harmonic\_with\_q.py}};

\node[kaldo] (DISPk) at (\colD,\rowKAL)
     {\textbf{kaldo}\\
      \texttt{phonons.frequency}\\
      \texttt{phonons.eigenvectors}\\
      \texttt{phonons.velocity}};

\node[kaldo] (FC3k) at (\colB,\rowKAL3)
     {\textbf{kaldo}\\
      \texttt{forceconstants.third}\\
      file: \texttt{THIRD} (binary)\\
      \texttt{third\_supercell=(3,3,3)}};

\node[rowlbl, color=kaldocol]
     at ([xshift=-0.4cm]POTk.west |- POTk) {kaldo};

% =========================================================
% PHONO(3)PY ROW
% =========================================================

\node[phono] (POTp) at (\colA,\rowPHO)
     {\textbf{phono3py (QE/PBE)}\\
      \texttt{dft(QE, PBE, Si.UPF)}\\
      \emph{opaque label}};

\node[phono] (FC2p) at (\colB,\rowPHO)
     {\textbf{phonopy}\\
      \texttt{phonopy --writefc}\\
      file: \texttt{fc2.hdf5}\\
      \texttt{phonopy.force\_constants}};

\node[phono] (DMp) at (\colC,\rowPHO)
     {\textbf{phonopy}\\
      \texttt{dynamical\_matrix}\\
      module:\\
      \texttt{harmonic/dynamical\_matrix.py}};

\node[phono] (DISPp) at (\colD,\rowPHO)
     {\textbf{phonopy}\\
      \texttt{get\_band\_structure\_dict()}\\
      \{frequencies, eigenvectors,\\group\_velocities\}};

\node[phono] (FC3p) at (\colB,\rowPHO3)
     {\textbf{phono3py}\\
      \texttt{phono3py-load}\\
      file: \texttt{fc3.hdf5}\\
      \texttt{Phono3py.fc3}};

\node[rowlbl, color=phonocol]
     at ([xshift=-0.4cm]POTp.west |- POTp) {phono(3)py};

% =========================================================
% MATERIALIZATION LINKS (abstract -> kaldo, kaldo -> phono)
% =========================================================

\foreach \abs/\kal/\pho in {POT/POTk/POTp, FC2/FC2k/FC2p, DM/DMk/DMp, DISP/DISPk/DISPp,
                            FC3/FC3k/FC3p}{
  \draw[matarr] (\abs.south)  -- (\kal.north);
  \draw[matarr] (\kal.south)  -- (\pho.north);
}

% =========================================================
% LEGEND (bottom strip)
% =========================================================

\begin{scope}[shift={(0, -10.0)}, font=\scriptsize]
  % small legend tiles
  \node[abstr, minimum width=1.6cm, minimum height=0.5cm] (lABS) at (0,0) {};
  \node[anchor=west] at (lABS.east) {\,abstract derived state};

  \node[param, minimum width=1.6cm, minimum height=0.5cm] (lPAR) at (5.3,0) {};
  \node[anchor=west] at (lPAR.east) {\,abstract parameter state};

  \node[kaldo, minimum width=1.6cm, minimum height=0.5cm, text width=, inner sep=2pt] (lK) at (10.6,0) {};
  \node[anchor=west] at (lK.east) {\,kaldo materialization};

  \node[phono, minimum width=1.6cm, minimum height=0.5cm, text width=, inner sep=2pt] (lP) at (15.0,0) {};
  \node[anchor=west] at (lP.east) {\,phonopy / phono3py materialization};

  % arrow legend (one row down)
  \draw[oparr]   (0, -0.7) -- ++(1.2, 0); \node[anchor=west] at (1.3, -0.7) {\,exact operator operation};
  \draw[stubarr] (5.3, -0.7) -- ++(1.2, 0); \node[anchor=west] at (6.6, -0.7) {\,stubbed (Phase 1) operation};
  \draw[matarr]  (10.6, -0.7) -- ++(1.2, 0); \node[anchor=west] at (11.9, -0.7) {\,materialization link};
\end{scope}

\end{tikzpicture}

\end{document}
